%=====================================================================
\chapter{The Publication System and Where Work Belongs}\label{ch:publication}
%=====================================================================
\locator{1}{Part~1 --- Foundations of Inquiry in Computing}

\begin{outcomes}
By the end of this chapter you will be able to:
\begin{tight}
  \item \textbf{Explain} why computing publishes in conferences when most
    disciplines do not, and what follows for your publication plan.
  \item \textbf{Interpret} the HJRS W, X, Y categories and state exactly what
    your degree requires of you.
  \item \textbf{Evaluate} a journal for legitimacy using three independent
    checks.
  \item \textbf{Critique} the impact factor and the h-index using the arguments
    the research-metrics community itself makes.
  \item \textbf{Describe} the peer review process from the reviewer's side, and
    why that view improves your writing.
\end{tight}
\end{outcomes}

\begin{prereq}
\chref{contribution}, for what a venue is being asked to accept. Read this
chapter early even though you will publish late: \reg{PhD Regs 2024}{\S15(iv)}
constrains \emph{when} your paper counts, and that constraint has to be planned
for, not discovered.
\end{prereq}

\begin{keyterms}
peer review $\bullet$ single- and double-blind review $\bullet$ rebuttal
$\bullet$ HJRS categories $\bullet$ impact factor $\bullet$ quartile $\bullet$
h-index $\bullet$ preprint $\bullet$ predatory journal $\bullet$ paper mill
$\bullet$ open access $\bullet$ gold, green and diamond $\bullet$ APC
$\bullet$ DOI $\bullet$ ORCID
\end{keyterms}

\section{Computing Publishes Differently}

In most sciences the conference is a place to present work in progress and the
journal is the archival record. In computing, for large parts of the field, it
is the other way round: the top conferences are selective, archival, and
carry more prestige than most journals~\cite{vrettas2015conferences}. ICSE,
CHI, NeurIPS, SIGCOMM and their peers run full peer review with acceptance
rates below 25\%, and their proceedings are the record.

This is a genuine disciplinary difference, and it creates a specific problem
for a scholar under Pakistani regulations, because the requirement in
\reg{PhD Regs 2024}{\S15} is written in terms of \emph{journals}.

\begin{regbox}[The publication requirement --- PhD Regulations 2024, \S15]
``For award of PhD degree, a PhD researcher shall be required to publish
research articles meeting the following criteria:
\begin{tight}
  \item[(a)] One research article in W category journal or two research
    articles in X category journals, for Science disciplines;
  \item[(b)] One research article in X category journal or two research
    articles in Y category journals, for Social Science disciplines.''
\end{tight}
And further: the researcher ``shall be the first author''; the article ``shall
be relevant to the PhD research work''; it ``shall be published \emph{after
approval of the research synopsis}''; and it ``shall be published in a relevant
research journal''.

\medskip
\textbf{There is no publication requirement for the MS IT degree}
(\cite{iub2024mphil}, \S47--48). MS scholars who publish do so for their own
reasons, which are good ones, but not to satisfy a regulation.
\end{regbox}

\begin{alertbox}[Three traps in \S15 that cost scholars a year]
\textbf{``After approval of the research synopsis.''} A paper accepted before
your synopsis is approved does not count towards the requirement, however good
it is. If you arrive with work in progress, sequence the synopsis first.

\medskip
\textbf{``First author.''} A paper on which your supervisor is first author
does not count for you. Agree authorship order in writing, at the start ---
\chref{integrity} gives the criteria to base that conversation on.

\medskip
\textbf{Conference papers are not journal articles.} A CORE~A conference paper
may be a far better contribution than a W-category journal article and still
not satisfy \S15. Plan for both: target the conference for the field, then a
substantially extended journal version for the regulation. Extension must be
genuine --- \chref{integrity} on text recycling and salami slicing.
\end{alertbox}

\section{Reading the Journal Categories}

The HEC Journal Recognition System~\cite{hec2024hjrs} assigns Pakistani-relevant
categories to journals, and these are what \S15 refers to. Categories are
reassigned periodically, so a journal that was X when you submitted may be
something else when you are examined.

\begin{center}
\small
\begin{tabular}{@{}L{1.6cm}L{5.4cm}L{7.5cm}@{}}
\toprule
\textbf{Category} & \textbf{Broadly} & \textbf{What it means for you} \\
\midrule
W & The strongest tier recognised
  & One paper here satisfies the Science-discipline requirement outright \\
X & Strong, widely indexed
  & Two papers here satisfy it; one satisfies the Social Science requirement \\
Y & Recognised
  & Two satisfy the Social Science requirement; insufficient alone for Science \\
Z / unlisted & Not recognised for this purpose
  & Does not count, whatever the journal claims about itself \\
\bottomrule
\end{tabular}
\end{center}

\begin{pitfall}[Check the category on the day you submit, and again on acceptance]
Categories move. Verify the journal's current HJRS status before submitting,
record a dated screenshot with your synopsis file, and check again at
acceptance. A scholar who submits to an X-category journal and is examined
after it has been downgraded has a problem that is entirely avoidable with two
minutes of record-keeping.
\end{pitfall}

\section{How Peer Review Actually Works}

Understanding review from the inside is the fastest way to improve your
writing, because it tells you what the reader is trying to do.

\begin{center}
\small
\begin{tabular}{@{}L{2.8cm}L{11.6cm}@{}}
\toprule
\textbf{Stage} & \textbf{What happens} \\
\midrule
Desk check     & An editor or chair rejects out-of-scope, unformatted or
                 obviously unsound submissions without review. A substantial
                 fraction of submissions die here, mostly avoidably \\
Assignment     & Two to four reviewers, chosen for topic fit. They are unpaid,
                 busy, and reading your paper late at night \\
Review         & Each reviewer forms a recommendation and writes it up.
                 Most decide provisionally within the first two pages \\
Discussion     & Reviewers see each other's reports and argue. This is where
                 a paper with one hostile review is often saved --- or lost \\
Rebuttal       & Many conferences let authors respond before the decision.
                 \chref{publishing} on how to write one \\
Decision       & Accept, minor or major revision, or reject. For journals,
                 revision cycles can repeat \\
\bottomrule
\end{tabular}
\end{center}

\begin{conceptbox}[What a reviewer is looking for, in order]
Reviewers are not looking for reasons to accept. They are looking, quickly, for
the answers to four questions --- and if a paper does not supply them in the
abstract and introduction, it starts from behind.

\begin{tightnum}
  \item \textbf{What is the claim?} If a reviewer cannot state your claim after
    the introduction, nothing later will rescue the paper.
  \item \textbf{Is it new?} Which is to say: does the related work section
    demonstrate you know what has been done, and locate your delta in it?
  \item \textbf{Is the evidence adequate to the claim?} The
    \chref{contribution} match between contribution type and evidence type.
  \item \textbf{Do the authors know the limits?} A candid threats section
    builds far more confidence than an absent one. Reviewers find the flaw
    regardless; the only question is whether you found it first.
\end{tightnum}
\end{conceptbox}

Blinding varies. \term{Double-blind} review hides author identity from
reviewers and is now common in software engineering and machine learning;
\term{single-blind} hides only the reviewer; \term{open} review publishes
reports and sometimes names. Each trades bias against accountability
differently, and none eliminates bias.

\section{Metrics, and Why the Metrics Community Distrusts Them}

You will be asked for your h-index. You will be told to publish in
``high-impact'' journals. It is worth knowing that the strongest criticisms of
these numbers come from bibliometricians, not from disgruntled authors.

\begin{center}
\small
\begin{tabular}{@{}L{2.9cm}L{4.6cm}L{6.9cm}@{}}
\toprule
\textbf{Metric} & \textbf{What it is} & \textbf{What it is not} \\
\midrule
Journal Impact Factor
  & Mean citations in year $t$ to items published in $t{-}1$ and $t{-}2$
  & A property of any individual paper. Citation distributions are heavily
    skewed: most papers in a high-IF journal are cited far below its IF \\
Quartile (Q1--Q4)
  & The journal's rank within a subject category
  & Comparable across categories, or stable --- category assignment changes
    the quartile \\
h-index
  & $h$ papers with at least $h$ citations each
  & Comparable across fields or career stages; it cannot decrease, and it
    rewards longevity over quality \\
Citation count
  & How often a work was cited
  & A measure of correctness or importance. Work is cited for being wrong,
    for being convenient, and for being famous \\
\bottomrule
\end{tabular}
\end{center}

\begin{paperbox}[The Leiden Manifesto]
Diana Hicks, Paul Wouters, Ludo Waltman, Sarah de Rijcke and Ismael Rafols,
\emph{Bibliometrics: The Leiden Manifesto for research metrics}, Nature 520,
2015~\cite{hicks2015leiden}.\oa{}

\medskip
\textbf{Why it matters.} Ten principles from the people who build these
indicators, on how they should and should not be used. Principle 1 ---
quantitative evaluation should \emph{support} qualitative expert assessment,
not replace it --- is the one most often violated, in Pakistan and everywhere
else.

\medskip
\textbf{What to extract.} The two principles that bear directly on you:
normalise for field before comparing anything, and never use journal-level
metrics to judge an individual article. Pair it with
DORA~\cite{dora2013declaration}, which commits signatories to the second.
\end{paperbox}

\begin{conceptbox}[Living with a system you have criticised]
There is an obvious tension: this section argues journal-level metrics should
not judge individual work, and \S15 requires you to publish in a
journal-category system. Both things are true, and the resolution is not
cynicism.

\medskip
Treat the category requirement as a \emph{floor} that must be cleared, and
venue \emph{fit} as the thing that actually determines whether your work is
read. A W-category journal whose readership has no interest in your subfield
will satisfy the regulation and be cited by nobody. A well-fitted X-category
journal in your area, plus a second X paper, satisfies the same regulation and
reaches the people whose work you want to influence.
\end{conceptbox}

\section{Predatory Venues}

A \term{predatory} journal charges a fee and provides no meaningful peer
review, while presenting itself as legitimate. The volume is
substantial~\cite{shen2015predatory}, the solicitation emails are relentless,
and graduate scholars under publication pressure are the target market.

\begin{alertbox}[The cost is not the fee]
Publishing in a predatory journal is worse than not publishing. The paper will
not satisfy \S15, because such venues are not HJRS-recognised. It cannot be
resubmitted elsewhere, because it is already published --- so the work is
destroyed, not merely wasted. And the line stays on your CV permanently, where
every future committee will see it.
\end{alertbox}

\subsection{Three independent checks}

Use all three. Any one alone can be fooled.

\begin{center}
\small
\begin{tabular}{@{}L{3.4cm}L{11.0cm}@{}}
\toprule
\textbf{Check} & \textbf{What to do} \\
\midrule
HJRS
  & Is the journal listed, and in which category? For your degree this is
    decisive regardless of anything else \\
DOAJ~\cite{doaj}
  & For open-access journals: is it indexed? DOAJ applies criteria on peer
    review, editorial process and transparency \\
Think.\ Check.\ Submit.~\cite{thinkchecksubmit}
  & Work its checklist. Do you recognise the editorial board? Is the peer
    review process described? Are fees stated clearly \emph{before}
    submission? Is the publisher a member of COPE? \\
\bottomrule
\end{tabular}
\end{center}

\begin{pitfall}[Warning signs, in rough order of reliability]
\begin{tight}
  \item An unsolicited email praising a paper of yours, and promising review in
    72 hours. Genuine journals do not work this way.
  \item A fee disclosed only after acceptance.
  \item An editorial board of people who cannot be found, or who did not know
    they were on it.
  \item Scope so broad it spans unrelated disciplines in one title.
  \item A claimed impact factor from an organisation nobody has heard of.
  \item Spelling errors on the journal's own home page.
\end{tight}
\medskip
Related, and increasingly serious: \term{paper mills} sell authorship on
fabricated manuscripts. Buying authorship is fabrication and falls under
\reg{PhD Regs 2024}{\S32} as academic misconduct, with the consequences
\chref{integrity} describes.
\end{pitfall}

\section{Preprints and Open Access}

A \term{preprint} is a manuscript posted publicly before or during peer review
--- arXiv is the computing norm. It establishes priority, makes work citable
immediately, and costs nothing. Most computing venues permit it; check the
specific venue, and note that a preprint is \emph{not} a publication for the
purposes of \S15.

\begin{center}
\small
\begin{tabular}{@{}L{2.3cm}L{5.0cm}L{6.9cm}@{}}
\toprule
\textbf{Route} & \textbf{How it works} & \textbf{What to watch} \\
\midrule
Gold    & Publisher makes it free; author or funder pays an APC
        & APCs of USD 2000--5000 are common and rarely fundable locally \\
Green   & Publisher version is paywalled; author deposits an accepted
          manuscript in a repository
        & Free. Check the embargo and which version is permitted \\
Diamond & Free to read \emph{and} free to publish; funded institutionally
        & Legitimate and growing; verify via DOAJ like any other \\
Hybrid  & A subscription journal offering paid open access per article
        & Widely criticised; pays twice for the same content \\
\bottomrule
\end{tabular}
\end{center}

\begin{conceptbox}[Green open access is the practical route]
For a scholar without a grant, green is almost always the answer: publish in
the best-fitting recognised journal, then deposit the accepted manuscript in an
institutional or subject repository. It costs nothing, is permitted by most
publishers after an embargo, and makes the work readable by people who cannot
pay. Register an \term{ORCID} identifier first, so that everything you deposit
accumulates against one unambiguous identity --- particularly valuable where
names are commonly shared.
\end{conceptbox}

\begin{traditions}{where work is published}
\tradEMP{Journals and archival conferences with full review. Increasingly
expects data and analysis scripts to be deposited alongside
(\chref{reproducibility}).}
\tradDSR{Information systems journals and design-oriented conferences.
Evaluation is what reviewers scrutinise, and the artefact itself is often
archived as a badged artefact.}
\tradQUAL{Journals that allow the length qualitative reporting requires; a
four-page limit cannot carry the thick description \reg{PhD Regs 2024}{\S14.3(A)}
asks for.}
\tradFORM{Theory conferences and journals where proofs are refereed in detail
and the archival version carries full appendices.}
\end{traditions}

\begin{lab}[Assess five venues]
Pick five journals or conferences plausibly relevant to your topic --- include
at least one that has emailed you an invitation.

\medskip
For each, record: HJRS category and the date you checked; DOAJ listing if open
access; whether it passes Think.\ Check.\ Submit.; the review model; typical
time to first decision; APC if any; and whether papers you have actually cited
appear there.

\medskip
That last column matters more than the others combined. A venue that has
published nothing you cite is probably not where your readers are. Rank the
five for your first paper and justify the ordering in three sentences.
\end{lab}

\begin{checkpoint}
\begin{tightnum}
  \item A PhD scholar in a Science discipline publishes one X-category paper
    before synopsis approval and one after. What does \S15 require of them now?
  \item Why does a high journal impact factor tell you little about a specific
    paper in that journal?
  \item Give two reasons a top-tier conference paper may be a better
    contribution than a W-category journal article, and one reason it may still
    not help you graduate.
  \item You are invited to submit to a journal that promises review within one
    week and lists an ``International Impact Factor'' of 6.2. State the three
    checks you would run and what each would likely show.
\end{tightnum}
\end{checkpoint}

\begin{chaptersummary}
\begin{tight}
  \item Computing treats top conferences as archival, but \S15 is written in
    terms of journals. Plan for both; a genuinely extended journal version of a
    conference paper is the standard route.
  \item PhD: one W or two X papers for Science, first author, after synopsis
    approval. MS: no publication requirement at all.
  \item Verify the HJRS category at submission and at acceptance, and keep
    dated evidence.
  \item Reviewers ask four questions in order: what is the claim, is it new, is
    the evidence adequate, do the authors know the limits. Write so the first
    two are answerable from the abstract.
  \item Journal-level metrics do not measure individual papers; the Leiden
    Manifesto and DORA say so from inside the field. Clear the category floor,
    then choose on fit.
  \item Predatory publication destroys work rather than wasting it. Check HJRS,
    DOAJ and Think.\ Check.\ Submit.\ --- all three.
  \item Green open access plus an ORCID is the practical, free route to being
    read.
\end{tight}
\end{chaptersummary}

\begin{reviewq}
  \item Computing's conference-centric culture is defended as fast and
    selective, and criticised as producing rushed, unreplicated work. Argue
    both sides, then state which you find more persuasive for your own
    subfield and why.
  \item \S15 requires publication before submission of the dissertation. Give
    the strongest argument for this requirement and the strongest argument
    against it, considering effects on both quality and on scholars' timelines.
  \item Double-blind review reduces some biases and introduces others. Name one
    of each, and propose a practical modification that would reduce the second
    without sacrificing the first.
  \item If journal-level metrics should not be used to evaluate individual
    researchers, what should be? Design an evaluation scheme for a hiring
    committee that is both fair and actually operable at scale, and identify
    what it would cost.
\end{reviewq}
